\textbf{Время:} 01:30

\textbf{Тезис:} три направления относительного импульса дают три характерных размера области испускания.

\textbf{На экране:} слева сначала ось пучка и long, затем средний поперечный импульс K\_T и out, затем side; справа условие LCMS и фраза о поперечном движении; в конце гауссова формула с тремя радиусами.

\textbf{Текст:} До сих пор мы обсуждали одну ширину корреляционного пика. Но область испускания не обязана иметь одинаковые размеры во всех направлениях. Поэтому воспользуемся тремя осями на рисунке.

Long направлено вдоль пучка. Out — вдоль среднего поперечного импульса выбранной пары, обозначенного K-T. Side перпендикулярно и пучку, и этому импульсу. Оси взаимно перпендикулярны; на плоской схеме они показаны в перспективе.

Для анализа нужно также выбрать систему отсчёта. Условие справа означает, что продольные импульсы двух частиц в сумме равны нулю. Такую систему называют LCMS. В неё переходим движением вдоль пучка, чтобы убрать общее продольное движение пары. Поперечное движение при этом остаётся: это ещё не система полного покоя пары.

Внизу та же гауссова форма, но теперь в показателе экспоненты три слагаемых. Для каждой компоненты q появляется свой радиус: out, side и long. Эти обозначения сохраняются на экспериментальных графиках, которые сейчас увидим.

\textbf{Переход:} Посмотрим, как три направления представлены в реальной корреляции STAR и почему её форма отличается от свободной гауссовой кривой.

\textbf{Справка:} LCMS — Longitudinally Co-Moving System. \(q^\mu=p_1^\mu-p_2^\mu\), \(K^\mu=(p_1^\mu+p_2^\mu)/2\); \(k_T=|\mathbf K_T|\) — модуль среднего поперечного импульса пары, не разность импульсов. Оси задаются для каждой пары, при k\_T=0 направление out не определено. Out не означает фиксированное направление наружу от центра источника. На слайде показано диагональное гауссово приближение без взаимодействия пионов после испускания; в общем случае необходимы смешанные члены. В анализе STAR смешанный член out–long может быть ненулевым из-за продольной асимметрии выборки; резерв содержит отдельный пример. Lisa et al. (2005), §2.4–2.5; Краева (2024), §2.1.
